% Introduction (owner SK). Every figure quoted here appears with its source in
% Section~\ref{sec:results}; nothing is stated here that is not measured there.
\section{Introduction}
\label{sec:intro}

Voice cloning has become a consumer technology. A few seconds of a person's speech
is enough for a zero-shot system to produce new sentences in their
voice~\cite{casanova2024xtts,betker2023tortoise}, and voice conversion does the same
without needing a transcript~\cite{rvc2023}. The obvious criminal application is the
telephone call: a voice the victim recognises, asking for money or a code, on a line
that carries no signal of who is really speaking.

Detection research has answered with strong results. On ASVspoof 2019
LA~\cite{todisco2019asvspoof}, self-supervised speech representations fine-tuned for
anti-spoofing reach error rates below one percent~\cite{tak2022wav2vec}. Our own
replication of that recipe reaches \textbf{0.85\% EER} on the evaluation partition.

That number is measured on English studio speech. The people most exposed to voice
fraud are not speaking English studio speech. In India, everyday conversation is
routinely \emph{code-mixed} --- Hindi and English alternating inside a single utterance
--- and it arrives over a narrowband telephone channel. This paper asks what happens
to a detector when both of those change at once, and what it costs to fix.

\subsection{The gap, and why it is not obvious}

The same checkpoint that scores 0.85\% EER on English scores \textbf{48.21\% EER on
clean code-mixed speech and 58.42\% over a simulated telephone channel}
(Section~\ref{sec:gap}). Over the channel its AUC is 0.381: below chance, meaning
the clones look \emph{more} genuine to it than the real speakers do.

The failure is not a mild degradation and it is not symmetric. At a 0.5 threshold
this model rejects 0.37\% of genuine English and 94.8\% of genuine Hindi--English ---
a 256-fold rise in the false-alarm rate on real speech. It has not become a weaker
detector of clones so much as a system that has stopped recognising real speakers.

Is it the code-mixing? A language-controlled test settles this. For 183 speakers
never used in training or testing, we take all-English, all-Hindi and mixed clips from
the \emph{same} lecture recordings, so only the language changes. The English-trained
detector rejects their genuine English as often as their mixed speech (84.3\% and
78.1\% on clean audio), and no amount of mixing, measured by the Code-Mixing Index,
predicts its score (Section~\ref{sec:res-language}). The gap is a recording-domain gap
that code-mixed telephone speech carries, not a property of the language mix, and
that is why adapting on in-domain audio can close it.

\subsection{Measuring against a floor, not against zero}

A low error rate on a code-mixed test set is easy to obtain for the wrong reason. If
the real and cloned clips differ in loudness, bandwidth, or recording room, a model
can separate them without learning anything about synthesis, and the resulting
number looks like detection.

We therefore report every result beside a \textbf{shortcut floor}: the error rate
eight cheap signal statistics reach on the same clips. A result better than its floor
is evidence that the model found something a trivial classifier could not; a result
at or above its floor is not evidence of anything. This discipline changes the
conclusions of this paper more than once --- including one of our own headline
conditions, which fails it.

\subsection{Contributions}

\begin{enumerate}
  \item \textbf{A measured gap.} The English-trained detector fails on code-mixed
    speech at 41.67--48.21\% EER clean and 44.15--58.42\% over the channel, far above
    every shortcut floor, with the failure characterised as a
    rejection of genuine speech rather than a loss of sensitivity to clones
    (Section~\ref{sec:gap}), and traced on a language-controlled set of 183 unseen
    speakers to the recording domain rather than to code-mixing
    (Section~\ref{sec:res-language}).
  \item \textbf{A cheap closure.} Low-rank adapters~\cite{hu2022lora} over a frozen
    wav2vec~2.0 encoder~\cite{baevski2020wav2vec}, training \textbf{1.13\% of the
    parameters}, bring channel-matched code-mixed error to \textbf{2.75\% EER}
    (3.27\% $\pm$ 1.17 over five training seeds) while retaining 2.42\% on clean
    English (Section~\ref{sec:systems}).
  \item \textbf{Matched shortcut floors.} Every detector is compared with a floor of
    eight cheap signal statistics fitted on its own training data and scored on its
    own clips, by a paired bootstrap over speakers. A floor fitted on the unseen tool
    itself would have understated the deployed detector (25.79\% against a matched
    37.79\% on Tortoise over the channel), and clip-level intervals would have
    overstated several of our own results (Section~\ref{sec:floor}).
  \item \textbf{An honest generalisation result.} Against a cloning tool no model was
    ever allowed to see, the adapters generalise on clean audio (5.44\% EER against a
    29.49\% floor). Over the telephone channel the deployed model averages 35.59\%
    across five training seeds against a 37.79\% floor and beats it in \emph{no} seed
    at the speaker level; full fine-tuning on a single attack family stays below its
    floor in every seed (19.72\%). Randomising the recording cues during training
    lowers the deployed recipe to 26.63\% over five seeds, which we report as
    exploratory (Section~\ref{sec:variants}). We report the weak result as
    prominently as the strong one.
  \item \textbf{Adaptation beats full fine-tuning where it matters.} Training the
    whole encoder on the same data reaches the lowest seen-tool error in the study
    (0.83\%) and collapses to chance on English. The adapters keep both.
  \item \textbf{A deployed system, not an illustration.} The detector runs on a live
    call and warns the callee --- and only the callee --- roughly ten seconds into
    the speech, using the same windowing and operating point the evaluation uses
    (Section~\ref{sec:system}).
\end{enumerate}

The corpora, attacks, manifests, configurations and per-clip scores behind every
number are committed, and each result table names the file it was read from.
